\section{Introduction}
Operating room (OR) and Operating theater (OT) planning and scheduling are central topics in healthcare operations research because the operating theater is often considered as one of the hospital's most resource-intensive and operationally critical environments. Efficient scheduling has direct implications for patient access, service quality, staff utilization, and hospital cost control. In literature, these planning problems are categorized in multiple decision levels, such as strategical, tactical and operational \cite{guerriero-2010}. \cite{CARDOEN2010921} reviews the literature based on problem characteristics. 

A key difficulty in OT scheduling is that surgery durations are inherently uncertain. Deterministic schedules based only on expected durations often perform poorly in practice, since deviations in realized procedure times propagate throughout the day and lead to patient waiting, downstream idle time, and overtime. Recent literature therefore increasingly models OR scheduling as an optimization problem under uncertainty.

In this research, an operational OT planning is made using mixed integer programming and sample average approximation. We study a weekly surgery scheduling problem for a set of inpatients that must be assigned to OT sessions, sequenced within each session, and given planned start times that are communicated in advance to the wards. The scheduling objective is to minimize a weighted sum of patient waiting time, OT idle time, and OT overtime, while respecting practical constraints such as fixed OT opening times, deterministic changeover times, and the operational requirement that patients cannot be brought to the OT before their announced start times. From a modeling perspective, this problem integrates assignment, sequencing, and timing decisions in a stochastic environment.

\section{Literature review}
The contribution of this report is to formulate the problem as a stochastic programming model and to evaluate its performance using a finite scenario set generated through sample average approximation (SAA). In this chapter, we will identify an appropriate solution approach and introduce the sample average approximation approach. 

\subsection{Solution approaches}
\subsubsection{Advance and allocation steps}
In literature, the advance and allocation steps can be solved using either sequential (SEQ) or integrated (INT) approaches \cite{Molina-Pariente2026}. Typically, the SEQ approach addresses the advance scheduling and allocation scheduling steps sequentially, whereas with the integrated approach, both advance and allocation steps are jointly addressed.

\textbf{Sequential Approach:}
Dai et al. propose a mathematical formulation to model surgery schedules under uncertainty of surgery duration in which patients are assigned to an operating room and place in the sequence of patients within the operation room \cite{DAI2025107098}.
In \cite{Denton2006} the one operation room problem is formulated as a two-stage stochastic mixed-integer-program, where the first state decision variables include the sequence of operations and planned starting time and the second stage decision variables include the idletime, overtime and waitingtime.

\textbf{Integrated Approach:}
\cite{Moosavi2018216} propose an integrated approach in which another index for the pre-defined number of slots available within a block is included in the decision variable such that a patient is assigned a day, session, and timeslot within the session. \cite{Vali-Siar2018549} uses a binary variable for whether a patient (p) is scheduled in OR (o) with surgeon (s) at time (t) on day (d). 

\subsubsection{Open (OS) and Block (BS) Scheduling}
Traditionally, open (OS) and block (BS) scheduling policies are considered in the literature. Under the OS policy, any surgeon can perform surgeries in any OR-day available within the planning horizon, which complicates the problem. A variant of the OS policy, the BS policy, is therefore frequently used in reality. With the BS policy, a surgeon only performs surgeries in one OR block time. In this case, the determination of the start times is simplified compared with the OS policy as it allows OR-sessions to be scheduled separately. 

To regard the advance and allocation steps, the sequential approach allows modelling the problem in a way that does not require pre-defining a total number of slots per session whhich is preferred given the problem at hand in which patient treatments durations differ significantly, such that it might require introducing dummy patients to fill all patient slots, depending on the model. In addition, the problem is dealt with as a block scheduling problem, in which blocks are separately sequenced to simplify the problem at hand. 

\subsection{Sample average approximation}
In the paper of Molina-Pariente et al. \cite{Molina-Pariente2026} an integrated operating room planning and scheduling problem under uncertainty is discussed. The model is formulated by denoting a set of days H (planning horizon) and a set J of ORs ($j=1,...,|J|)$ available on each day h in the planning horizon, denoting (j,h) as OR-day. To solve the problem, the sample average approximation (SAA) procedure, in which the problem is approximated as a discrete deterministic problem that integrates N samples, is proposed. \cite{Denton2006} also proposes SAA to solve the problem at hand. The SAA approach \cite{Kleywegt} follows the logic described in algorithm \ref{alg:saa} and is a good solution methodology for problems in which numerous different outcomes may occur, such that an exact solution becomes computationally intractable. \textbf{For such problems, stochastic programming using SAA is a good way to balance computational complexity and the solution quality by considering various scenarios all at once, as if it were deterministic. The disadvantage is that the solution quality solely depends on how well the used scenarios represenet the potential outcomes. In the SAA design, a proper sampling method as well as an appropiate sample size (which balances computational requirements and outcome coverage) should be determined for the matematical model. In addition, an appropiate evaluation sample should be determined to evaluate the solution quality upon. ~"show why SP with SAA is a good solution methodology for this problem" "From this review it should be clear why SAA is suitable for the problems at hand, what the disadvantages are, which design choices should be made, and a step-by-step approach on how you will implement SAA in your assignment."}

\begin{algorithm}[H]
\caption{Sample Average Approximation (SAA) Algorithm for Stochastic Discrete Optimization}
\label{alg:saa}
\begin{algorithmic}[1:]
\Statex \textbf{Step 1:}
\State Choose initial sample sizes $N$ (nr. of scenarios) and $N'$ (size of the evaluation sample), a decision rule for determining the number $M$ of SAA replications, a decision rule for increasing the sample sizes $N$ and $N'$ if needed, and tolerance $\varepsilon$.
\Statex \textbf{Step 2:}
\For{$m = 1, \dots, M$}
    \State Generate a sample of size $N$ and solve the SAA problem with objective value $\hat{v}_N^m$ and $\varepsilon$-optimal solution $\hat{x}_N^m$.
    \State Estimate the optimality gap $g(\hat{x}_N^m) - v^*$ and the variance of the gap estimator.
    \If{the optimality gap and the variance of the gap estimator are sufficiently small}
        \State \textbf{go to Step 4}
    \EndIf
\EndFor
\Statex \textbf{Step 3:}
\If{the optimality gap or the variance of the gap estimator is too large}
    \State Increase the sample sizes $N$ and/or $N'$.
    \State Return to Step 2.
\EndIf
\Statex \textbf{Step 4:}
\State Choose the best solution $\hat{x}$ among all candidate solutions $\hat{x}_N^m$ produced, using a screening and selection procedure.
\State \textbf{Stop.}
\end{algorithmic}
\end{algorithm}

\section{Problem description}

\section{Mathematical Model}

\section{SAA procedure}

\section{Results and Evaluation}